\section{The Pentacomb Resolution}

The thesis of this paper places the physics on the $\{3,4,3,4\}$ mesh. This
section is not that. It is a selection study, a careful look at one near-relative
of $\{3,4,3,4\}$ to see what a substrate that carries both spin and curvature
actually looks like. The substrate is the $5$D pentacomb $\{3,4,3,3,4\}$. It is
studied here for comparison, to close one open trade in the selection argument, and
then set aside. The living physics stays on $\{3,4,3,4\}$.

The trade in question is sharp. Among the substrates examined so far, one carries
the spinor directions but is flat, and another is curved but its directions carry
no spinor. It looked as though one had to give up exactly one of the two to get the
other. The pentacomb shows that one does not. It has the spinor directions
\emph{and} the curvature, and a fermion propagates on it from the bare directional
knit, with no added field and no change to the law. Spin and curvature need not be
traded.

Two terms are used throughout, so a plain statement of each helps. A \emph{site} is
one of the $24$ directions out of a dock, the directed slot a tone lives in. The
\emph{knit} is the law of the mesh, collide then stream, a fixed reversible
charge-conserving table run once per beat. A \emph{spinor} is the object that
describes half-integer spin. Its mark is that it must be turned twice around, through
$720$ degrees, before it returns to itself. The question of this section is whether
the directions of the sites can themselves carry that double-cover sign while the
mesh stays curved.

\begin{result}[The decisive statement]
The spin-versus-curvature trade is resolved by the substrate that has both the
directions and the dimension. The $5$D pentacomb $\{3,4,3,3,4\}$ carries the $D_4$
spinor directions and a negatively curved metric at once, and a fermion propagates
on its bare directional knit.
\end{result}

This is route four of the four routes laid out in the spin obstruction section. It
is the only one of the four that gets a propagating spinor from the bare directional
knit. Its depth is $L2$, a measured propagation result, not the full $L3$ continuum
claim.

\subsection{The trade, stated}

Two well-studied substrates each fail one half of what is wanted. A useful matter
substrate would carry the spinor in its very directions, so that no extra field is
needed, and it would be negatively curved, so that it can grow forever and host a
boundary. The flat $24$-cell honeycomb has the first property and not the second.
The hyperbolic $\{5,3,4\}$ has the second and not the first.

\begin{table}[ht]
\centering
\begin{tabular}{@{}llll@{}}
\toprule
substrate & curvature & spinor in the sites & propagating fermion from bare knit \\
\midrule
flat $\{3,4,3,4\}$ & none (flat) & yes, the $24$ $D_4$ directions & yes, but flat \\
hyperbolic $\{5,3,4\}$ & negative & no, the $12$ directions permute & no, needs a detour \\
pentacomb $\{3,4,3,3,4\}$ & negative & yes, the $24$ $D_4$ directions & yes \\
\bottomrule
\end{tabular}
\caption{The spin-versus-curvature trade and its resolution. The flat $24$-cell
honeycomb has the spinor directions but no curvature. The hyperbolic $\{5,3,4\}$ has
curvature but its $12$-direction sites carry no spinor. The pentacomb has both.}
\label{tab:pentacomb-trade}
\end{table}

The $\{5,3,4\}$ failure is a theorem, not a difficulty of measurement. Its $12$
directions form the icosahedral group, and a permutation of directions is a linear
action. A spinor is projective, it carries the double-cover sign, and a linear action
cannot carry that sign. So the bare directional knit on $\{5,3,4\}$ has no spinor in
its sites. Matter can still live on $\{5,3,4\}$, but only through the form complex or
through a tone defect, a detour. Each of the two known substrates lacks exactly what
the other has. The pentacomb lacks neither.

\subsection{Why the pentacomb has the spinor directions}

The pentacomb $\{3,4,3,3,4\}$ contains the $24$-cell as a substructure. The $24$
directions of the $24$-cell are the $D_4$ roots. Under the symmetry of those roots
the $24$ directions split into three eight-dimensional sectors.

\begin{table}[ht]
\centering
\begin{tabular}{@{}lll@{}}
\toprule
sector & what it is & carries the spinor \\
\midrule
$8v$ & vector & no \\
$8s$ & spinor & yes, the minus-one \\
$8c$ & conjugate spinor & yes, the minus-one \\
\bottomrule
\end{tabular}
\caption{The $D_4$ directions split as $8v + 8s + 8c$. The two spinor sectors carry
the double-cover sign that the linear permutation field of $\{5,3,4\}$ cannot.}
\label{tab:pentacomb-triality}
\end{table}

The $8s$ and $8c$ are genuine spinor sectors. They carry the minus-one under a $2\pi$
rotation. The order-three symmetry that exchanges the three eights is triality, the
outer symmetry of $D_4$. The point is that the directions of the sites are themselves
a spinor-carrying object. No detour through a form complex or a defect is needed. The
double-cover sign sits in the sites from the start.

\begin{proposition}[The spinor sites of the pentacomb]
The pentacomb contains the $24$-cell, whose $24$ directions are the $D_4$ roots.
Those roots decompose as $8v + 8s + 8c$, with $8s$ and $8c$ genuine half-integer-spin
sectors. The directional sites of the pentacomb therefore carry the spinor directly,
not through any added index. Triality permutes the three eights.
\end{proposition}

\subsection{Why the pentacomb is curved}

The flat $24$-cell honeycomb $\{3,4,3,3\}$ already has these same spinor directions,
yet it is flat. To get the curvature one needs a high enough dimension to be
hyperbolic while still containing the $24$-cell. The pentacomb meets this at five
dimensions. Its Coxeter Gram matrix, the matrix that records the angles between the
mirrors that generate the tessellation, is Lorentzian. A Lorentzian Gram matrix is
the signature of negative curvature, exactly as a positive-definite one is the
signature of flat space.

\begin{table}[ht]
\centering
\begin{tabular}{@{}ll@{}}
\toprule
object & property \\
\midrule
$\{3,4,3,3,4\}$ Coxeter Gram matrix & Lorentzian (negative curvature) \\
$\{3,4,3,3\}$ Euclidean $24$-cell honeycomb & flat (the control) \\
\bottomrule
\end{tabular}
\caption{The pentacomb is hyperbolic by its Lorentzian Gram matrix. The flat
$24$-cell honeycomb is its control, same directions, no curvature.}
\label{tab:pentacomb-curvature}
\end{table}

The flat $24$-cell honeycomb is the control that isolates curvature as the variable.
It has the same $24$-cell and the same $D_4$ directions, but it has no curvature. If
a result holds on the pentacomb and not on this flat control, the curvature is doing
the work. If it holds on both, the result belongs to the directions, not the
curvature.

\subsection{A fermion propagates}

The structural facts above say the spinor is present and the curvature is present.
The dynamical question is whether a field can actually move. This is settled by the
pentacomb propagation experiment, depth $L2$. It runs the same return-probability
measure used for matter on $\{5,3,4\}$, now on the generated pentacomb cell graph.

The measure is the quantum return probability, how much of a localized excitation
stays at its starting dock under the unitary evolution of the operator. A low return
means the excitation spreads away and does not come back, the extended propagating
phase. A high return means it is trapped, the localized phase. The control that
makes the propagating reading meaningful is a strong deterministic disorder
potential, no randomness, which should localize the excitation and drive the return
back up. The base is discrete and deterministic throughout. No randomness enters.

\begin{table}[ht]
\centering
\begin{tabular}{@{}lll@{}}
\toprule
run & return probability & phase \\
\midrule
pentacomb $\{3,4,3,3,4\}$, clean & about $0.05$ & extended, propagating \\
strong deterministic disorder & about $0.5$ & localized \\
flat $24$-cell honeycomb, clean (control) & propagates & extended \\
\bottomrule
\end{tabular}
\caption{A Kahler-Dirac fermion propagates on the pentacomb. The clean return is
about $0.05$ in the extended phase, against about $0.5$ under deterministic disorder.
The flat control propagates too, confirming a spinor on a curved substrate, not
merely the flat $24$-cell. The result is robust across two sizes.}
\label{tab:pentacomb-propagation}
\end{table}

A Kahler-Dirac fermion propagates on the pentacomb. The clean return is about $0.05$
in the extended phase, against about $0.5$ when the deterministic disorder is turned
on. The fermion localizes when it should and propagates when it should. This is a
control that could have failed, since a fermion already pinned by the hyperbolic
irregularity would have returned high on the clean run too. It does not. The result
is robust across two lattice sizes.

The flat $24$-cell honeycomb propagates as well. That is the discriminating control.
It confirms that what propagates on the pentacomb is a spinor on a \emph{curved}
substrate and not merely the flat $24$-cell wearing a different name.

\begin{result}[The trade is resolved]
The pentacomb has the spinor sites, the curvature, and a propagating fermion, all
three together. Flat $\{3,4,3,4\}$ has spin but no curvature. Hyperbolic $\{5,3,4\}$
has curvature but spinless sites. The pentacomb lacks none of the three.
\end{result}

\subsection{The bare knit runs on it}

A spinor on a curved metric is not yet the theory's object unless the bare knit, the
same reversible charge-conserving collide-then-stream law, actually runs on the real
pentacomb mesh. It does. The generated $5$D pentacomb cell graph, built from
thousands of docks by its Coxeter reflections, grows faster than the flat honeycomb,
as a negatively curved space should. The reversible charge-conserving knit runs
directly on its real adjacency, conserving charge and reversing exactly.

So the spinor comes from the directions of the sites, not from any added index, and
the law is unchanged. One extra spatial dimension buys the spinor sites for free.

\begin{principle}[No change to the law]
The pentacomb resolution adds no field and changes no part of the knit. The spinor is
read off the directional sites, and the same reversible charge-conserving knit runs
on the real pentacomb adjacency. The whole cost is one extra spatial dimension.
\end{principle}

\subsection{Why this is the decisive resolution}

The spin obstruction section gave four routes around the obstruction. It is worth
setting them side by side, because only one route gets a propagating spinor from the
bare directional knit with nothing added.

\begin{table}[ht]
\centering
\begin{tabular}{@{}lll@{}}
\toprule
route & what it changes & propagating spinor, no added field \\
\midrule
1. defect & nothing & no, a heavy soliton on an unsolved persistence problem \\
2. spinor index plus connection & the knit & no, it adds a field \\
3. staggered Kahler-Dirac & the knit & no, it adds the form structure \\
4. pentacomb & the substrate & yes \\
\bottomrule
\end{tabular}
\caption{The four routes around the spin obstruction. Routes one through three each
add a field or ride on an unsolved problem. Route four, the pentacomb, changes only
the substrate.}
\label{tab:pentacomb-routes}
\end{table}

Routes one through three either add a field or ride on an unsolved problem. The
defect route adds nothing to the law but yields a heavy soliton that rests on the
unsolved persistence problem. The index-plus-connection route and the staggered
Kahler-Dirac route both change the knit, one by adding a spinor field, the other by
adding the form structure. Only route four changes nothing but the substrate.

\begin{result}[Route four is unique]
Route four is the only way to get a propagating spinor from the bare directional
knit, with no added field. The cost is one extra spatial dimension. The benefit is
that spin and curvature stop being a trade. Both are intrinsic to the same mesh, spin
in the directions of its sites and curvature in its Gram matrix.
\end{result}

\subsection{The honest scope}

This result is $L2$, not $L3$, and the boundary is stated plainly.

What is established is propagation. The extended phase is confirmed by the return
probability against a deterministic-disorder control, on a discrete and deterministic
base, with no randomness.

What is not claimed is the full relativistic Dirac-cone dispersion. The irregularity
of the mesh makes a clean momentum-space dispersion hard to read, the same residual
that remains for matter on $\{5,3,4\}$. That refinement stays open.

What backs the cell-graph knit is real geometry. The full $5$D hyperbolic embedding
is built, and greedy routing on $H^5$ delivers, so the knit that runs on the
combinatorial graph is backed by an actual hyperbolic embedding and not by a
combinatorial accident. What is settled is the trade. What is open is the continuum
dispersion.

\subsection{Index of verified results}

\begin{table}[ht]
\centering
\begin{tabular}{@{}ll@{}}
\toprule
claim & depth \\
\midrule
a fermion propagates on the $5$D pentacomb (spin sites and curvature, trade resolved) & $L2$ \\
the pentacomb is hyperbolic and carries the $D_4$ spinor directions ($8v + 8s + 8c$) & $L2$ \\
the bare knit runs on the generated pentacomb mesh & $L2$ \\
the full $5$D hyperbolic embedding, greedy routing on $H^5$ delivers & $L2$ \\
\bottomrule
\end{tabular}
\caption{The verified results of the pentacomb resolution, all at depth $L2$.}
\label{tab:pentacomb-index}
\end{table}

This section closes one open trade in the selection argument and then steps back. It
shows that a substrate can carry the spinor in its directions and curvature in its
metric at the same time, and that a fermion moves on it from the bare knit. The
physics of the theory still lives on $\{3,4,3,4\}$. The pentacomb is its
spinor-carrying curved near-relative, studied so the comparison line is complete.
